\subsection{Harness 群体与评估设置}
\label{sec:interface}
\emph{构建器}编写 Python harness，\emph{目标求解器}回答程序发出的模型调用，外部\emph{评分器}判定最终输出。所报告的比较使用 \texttt{GLM-5.3-Flash}（附录~\ref{app:models}）。每个 harness 实现 \texttt{solve(question)}，通过 \texttt{self.llm} 调用模型；BIRD 还通过 \texttt{self.execute} 提供数据库结构与 SQL 反馈。标准答案仅对评分器可见。\bare{} 基线执行一次温度为 $0$ 的调用。生成程序可以改变提示、解码、采样次数与控制流；这些设置决定各成员的计算过程与采样行为。表~\ref{tab:study-map} 说明各项研究的任务。

\begin{table}[t]
\centering\small
\setlength{\tabcolsep}{4pt}
\begin{tabular}{@{}p{0.21\linewidth}p{0.29\linewidth}p{0.19\linewidth}p{0.23\linewidth}@{}}
\toprule
研究 & 程序 & 评估 & 回答的问题 \\
\midrule
MATH 探索 & 35 个生成程序 + \bare{} & $\RcvPairedN$ 题；单次运行 & 是否存在错误覆盖？ \\
MATH 重复对照 & 8 个生成程序 + \bare{}，对照 9 个克隆 & 同一批 $\RcvPairedN$ 题；3 次新重复 & 选择能否跨重复迁移？ \\
BIRD 发现 & 12 个生成程序 + \bare{} + \react{} & 60 题；单次运行 & 什么机制实际执行？ \\
BIRD 准入 & 6 个构建器 $\times$ 3 个种子 $\times$ 4 个臂 & 主比较 1,169 题；全臂 400 题 & 筛选是否改变覆盖？ \\
\bottomrule
\end{tabular}
\caption{\textbf{每项研究支持不同的判断。}单次运行指每个成员—任务对执行一次。MATH 重复面板从探索性群体中抽取。BIRD 提供独立的执行与准入证据，各项研究分别报告效应。}
\label{tab:study-map}
\end{table}

\paragraph{MATH 群体与面板抽取。}
\label{sec:math-setup}
六个构建器为 GLM-5.3、Qwen3.8-Max、DeepSeek-V4-Pro、Kimi-K3、MiniMax-M3 和 ERNIE-5.0-Thinking-Preview，每个在一个种子下获得八个槽位。每个槽位是独立的生成机会，并非指定策略。提示要求改进单次贪心调用，程序结构由构建器自行设计，不施加跨槽位多样性约束。重试复用相同提示，不反馈前次代码或评估结果。每槽最多尝试三次（温度 $0.7$，每次最多 $32{,}768$ 个输出 token），在首个程序通过解析、导入、接口及两道开发题的非空答案检查后停止。$48$ 个槽位得到 $35$ 个冻结程序，用于后续评估。附录~\ref{app:builder-instruction} 概述构建器指令和接口约束。

重复面板按历史平均调用次数不超过 $3.05$ 的条件保留 $21$ 个合格成员，再由固定加盐哈希选择每个构建器一个席位和两个自由席位。这偏向低调用程序。所选源码包含起草—检查—修订、一致性检查和采样投票；逐项清单见附录~\ref{app:panel-mechanisms}。加入 \bare{} 后共九个成员。该面板与九个相同基线槽均进行三次新执行，关闭缓存，解码设置由各成员源码固定。

本文使用 MATH-500 \citep{math500} 中的 $100$ 道开发题与同一批 $\RcvPairedN$ 道评估分析题。该分析集在采集后按完整观测条件确定：两臂全部重复和单次评估群体均有评分；$\RcvExcludedN$ 道重复观测不完整的题被排除。本文所有 MATH 主比较使用相同题目身份，分析集内评分单元观测完整率为 $100\%$。重复研究在这些任务上进行新执行。排除情况见附录~\ref{app:wp1r}。
评分器抽取最终答案，优先使用所要求的 \texttt{\#\#\#\#} 行；两端均可解析为数值时比较数值，否则比较规范化 LaTeX。比较器及其审核详见附录~\ref{app:crossdomain}。

\paragraph{BIRD 发现与准入。}
\label{sec:bird-setup}
BIRD \citep{bird2021} 要求根据问题与数据库生成 SQL。TTHE 的 proposer 使用 GLM-5.3-Flash 构建器产生 $12$ 个发现候选，经导入与冒烟执行检查准入，与两个自带基线在 $60$ 题上评估。独立的准入研究交叉自由／指定策略与无门控／合规门控，使用六个构建器、三个种子、每臂八个槽位。每槽最多三次尝试（温度 $0.7$，$16{,}384$ 个输出 token）。指定策略包括执行修复、三次采样投票及 schema linking。所有臂均要求接口有效、SQL 非空；门控臂额外检查机制合规，并排除部分结构。

准入研究保留两个数据库的 $365$ 题用于开发，九个数据库的 $1{,}169$ 题用于评估；四臂共用预先抽取的 $400$ 题核心集。一次 $18$ 题冒烟测试是已披露的开发暴露例外。预测与标准查询的结果由官方行集合评分器比较，超时为 $30$ 秒，执行失败计零。发现研究则使用较早的字符串化行代理，并限制返回行数。准入采集使用共享响应缓存，MATH 重复执行则关闭缓存。这些条件及门控干预的捆绑性质保留于附录~\ref{app:phase1-provenance} 和~\ref{app:phase2-protocol}。
